\documentclass[10pt,twocolumn]{article}

% Packages for CVPR/ICCV/ECCV format
\usepackage[utf8]{inputenc}
\usepackage[T1]{fontenc}
\usepackage{amsmath,amsfonts,amssymb}
\usepackage{graphicx}
\usepackage{xcolor}
\usepackage{hyperref}
\usepackage[numbers,compress]{natbib}
\usepackage{geometry}
\geometry{margin=0.8in}

% Title and author information
\title{Spatiotemporal Object Detection: Fusing RGB Images with Direct Time-of-Flight Transients for Enhanced Material-Aware Scene Understanding}
\author{A. Akuoko, D. Chitnis, and I. Gyongy}
\date{\today}

\begin{document}

\maketitle

\begin{abstract}
This paper presents the first reported combination of two-dimensional spatially resolved RGB images with one-dimensional direct time-of-flight (dToF) time-resolved signals for enhanced object detection. Unlike RGB-D systems that estimate depth via stereoscopy or indirect time-of-flight, this solution employs dToF signals (photon timing histograms) that provide more than depth: they capture optical properties useful for inferring material composition. A hybrid neural network architecture combines time- and space-resolved signals, with object localisation and label features derived from spatially resolved images and structural features extracted from time-resolved signals. Experimental results demonstrate that even with narrow-band or single-wavelength achromatic signals, machines achieve improved structural and spatial understanding when structural or depth features originate from transient signals and topological features from spatially resolved images. We evaluate three state-of-the-art spatial object detection models (YOLOv3, YOLOv8, and DINOv3) integrated with a custom material classification model, achieving combined detection performance of 94.5\%--98.5\%, with YOLOv8 achieving the highest combined performance (98.5\%) at near-real-time processing speeds ($\sim$8\,ms). The material classifier achieves exceptional performance (99\% validation accuracy) while operating at extremely fast inference speeds (1--5\,ms), demonstrating that structural features extracted from time-resolved signals significantly improve machine vision quality and reliability.
\end{abstract}

\section{Introduction}

Automatic pattern recognition systems for machine vision have achieved unprecedented levels of sophistication, yet contemporary computing systems remain substantially inferior to the visual perception capabilities exhibited by advanced biological systems. A fundamental limitation underlying numerous perception quality challenges is the exclusive reliance on two-dimensional spatially resolved images for extracting structural features that are more accurately inferred from time-resolved signals.

This work presents an instrument designed to enhance structural feature extraction in machine vision systems through the integration of a single-photon avalanche diode (SPAD) time-of-flight sensor with an active pixel sensor (APS). The experiments demonstrate that, even in the absence of spectral composition analyses and without reliance on high computational power or sophisticated neural network architectures, the combination of transient and spatially resolved intensity images of the same target enables machines to perceive scenes in a manner that enhances higher-order signal intuition, thereby improving reasoning capabilities compared to contemporary vision systems.

The main experimental input combines two-dimensional intensity images with one-dimensional transient images. When these inputs are evaluated using state-of-the-art (SOTA) visual perception models, improvements in machine intuition are observed. Machines demonstrate enhanced understanding of visual content with respect to true structural details, which shapes their intuition for making more reliable probabilistic decisions.

\section{Related Work}

State-of-the-art vision research demonstrates impressive work in three-dimensional perception and reconstruction, a task strongly founded on the quality of structural features for depth estimation \cite{laga2022}. Nevertheless, the majority of recognised research in this domain has implemented solutions by depending on spatially resolved signals for all features. Scene flow prediction, three-dimensional image attribute estimation, three-dimensional reconstruction with diffusion models, and related tasks continue to advance rapidly along computational performance benchmarks while remaining susceptible to adversarial examples.

Modern machine-vision pipelines that fuse RGB imagery with co-registered dToF transients demonstrate improved robustness where spatial appearance is ambiguous or adversarially spoofed \cite{moramartin2021,yang2024}. However, most existing approaches focus on depth estimation rather than material-aware detection. This work extends prior fusion approaches by explicitly leveraging temporal signals for material classification while maintaining object detection capabilities.

\section{Methodology}

\subsection{System Architecture}

The proposed system employs a dual-model architecture comprising:
\begin{enumerate}
    \item \textbf{Spatial detection module}: Processes two-dimensional RGB images for object localisation and classification
    \item \textbf{Material detection module}: Processes one-dimensional transient signals for material composition classification
\end{enumerate}

The spatial detection module accepts 224$\times$224 pixel RGB images and outputs bounding box coordinates, object class labels, and confidence scores. The material detection module accepts processed 16$\times$16 pixel transient images and outputs material class probabilities.

\subsection{Spatial Detection Models}

Three state-of-the-art spatial object detection models were integrated:

\textbf{YOLOv3}: A pre-trained model comprising approximately 26 million parameters, fine-tuned on a custom dataset of eight object classes \cite{redmon2018}. The model achieved validation accuracy in the range of 94\%--99\%, with inference time of approximately 99\,ms for 224$\times$224 pixel inputs on CPU hardware.

\textbf{YOLOv8}: A more recent iteration employing a Cross Stage Partial (CSP) style block architecture. The model achieved mean Average Precision (mAP) of approximately 96.3\%, with inference speed approximately 1.5$\times$ faster than YOLOv3 (68\,ms for 224$\times$224 pixel inputs).

\textbf{DINOv3}: A self-supervised vision transformer framework employing a teacher-student knowledge distillation paradigm \cite{simeoni2025}. The model achieved approximately 97\% accuracy, demonstrating the effectiveness of self-supervised learning for object detection tasks.

\subsection{Material Detection Module}

The material detection module employs a custom convolutional neural network architecture optimised for material discrimination from transient image inputs. The classifier accepts processed 16$\times$16 pixel images, converted to tensor format and normalised prior to feature extraction.

The network architecture comprises a three-block convolutional structure with progressive feature channels (1, 3, 32, 64, 128), enabling hierarchical feature extraction from low-level edge detection to high-level material-specific patterns. The architecture utilises global average pooling to reduce spatial dimensions, mitigating overfitting while preserving discriminative feature representations.

Three distinct model variants were trained according to different class consolidation strategies, yielding 34-class, 18-class, and 12-class configurations. Performance evaluation revealed that validation accuracy approached 99\% for consolidated class sets, whereas the full 34-class configuration achieved approximately 60\% accuracy.

\section{Experimental Setup}

\subsection{Dataset}

The experimental dataset comprises eight object classes: \texttt{bell\_pepper}, \texttt{bowl}, \texttt{carrot}, \texttt{eggplant}, \texttt{poor\_vis}, \texttt{potato}, \texttt{teacup}, and \texttt{tomato}. The \texttt{poor\_vis} label represents images acquired under low-light conditions, demonstrating that the material detection module can successfully identify materials from corresponding transient images captured under identical low-light conditions.

Training was conducted on 640 images with 160 validation samples. Both training and inference procedures were configured to execute on CPU hardware to demonstrate practical deployment feasibility.

\subsection{Signal Processing}

Temporal signals from SPAD sensors are processed using the per-bin Gaussian approximation model described in the PAWD framework \cite{tontini2020,koerner2021}. For multi-material surfaces, the expected photon count per bin is derived from a Gaussian mixture model:

\begin{equation}
\mu_i = \sum_{m=1}^{M} \frac{S_m}{\sqrt{2\pi\sigma_m^2}} \exp\left(-\frac{(t_i - t_{0,m})^2}{2\sigma_m^2}\right) + b
\end{equation}

where $S_m$ is the total signal photons for material $m$ across the entire peak, $\sigma_m$ is the standard deviation, $t_{0,m}$ is the peak time for material $m$, and $b$ represents ambient noise per bin.

\section{Results}

\subsection{Individual Model Performance}

YOLOv8 achieved the highest object detection accuracy at approximately 98\%. DINOv3 achieved approximately 95\% accuracy, demonstrating the effectiveness of self-supervised learning. YOLOv3 achieved approximately 94\% validation accuracy. The material classifier achieved exceptional performance (99\% validation accuracy) while operating at extremely fast inference speeds (1--5\,ms).

\subsection{Combined Dual Detection Performance}

Table~\ref{tab:combined_performance} summarises the performance per spatial model and the corresponding performance in combination with the material classifier. The combined performance is calculated as a simple average, assuming equal relevance for both material and object classification correctness.

\begin{table}[h]
\centering
\caption{Combined dual detection performance}
\label{tab:combined_performance}
\begin{tabular}{lcc}
\hline
Model & Spatial Accuracy & Combined Accuracy \\
\hline
YOLOv3 & 94\% & 94.5\% \\
YOLOv8 & 98\% & 98.5\% \\
DINOv3 & 95\% & 95.5\% \\
\hline
\end{tabular}
\end{table}

YOLOv8 combined performance and speed are approximately 98.5\% and $\sim$8\,ms, respectively, demonstrating that the integration maintains high accuracy while providing dual detection capabilities suitable for real-time applications.

\section{Discussion}

The high validation accuracies achieved across multiple independent model architectures provide evidence that the spatiotemporal approach learns generalisable representations rather than overfitting to training data. The consistency of high performance across architecturally distinct models indicates that the observed accuracy stems from the complementary nature of spatial and temporal features rather than model-specific memorization.

The material classifier's exceptional performance is particularly justified given the well-characterised resultant signal profile described by the Gaussian mixture model. The material classifier's ability to achieve high accuracy while operating at extremely fast inference speeds demonstrates that the model has learned efficient, generalisable feature representations from well-modelled transient signal characteristics.

\section{Conclusion}

This work demonstrates, for the first time, the successful fusion of RGB spatial images with dToF temporal signals for enhanced object detection with material awareness. The results show that structural features extracted from time-resolved signals significantly improve machine vision quality, enabling more reliable scene understanding compared to spatial-only approaches. The system achieves near-real-time performance (8\,ms) while maintaining high accuracy (98.5\%), making it suitable for practical deployment in safety-critical applications.

\bibliographystyle{IEEEtranN}
\bibliography{../references}

\end{document}

